% ============================================================================
% Section: When Does Adaptive Safety Pay Off?
% ============================================================================
% Purpose: Characterize the cost-benefit frontier across prior quality
% regimes using the prior degradation sweep.
% All Corralling results use w₀=0.7 (recommended prior-trust bias).

\subsection{When Does Adaptive Safety Pay Off?}
\label{sec:metalearning_cost}

Having established that semantic priors enable strong routing
performance (\S\ref{sec:results}), we now ask a sharper question:
\textbf{when should a practitioner use each of the library's three
routing strategies?}

The answer---a crossover experiment spanning the full spectrum from
correct to adversarial priors---reveals a crossover point that
partitions prior quality into two regimes, with Corralling providing
bounded regret across both.

% ----------------------------------------------------------------------------
% Experimental Setup
% ----------------------------------------------------------------------------
\paragraph{Experimental Setup.}
We evaluate three routing strategies across 11~prior corruption levels
on the 750-prompt LMSYS holdout set (2~models, moderate distribution
shift).  Prior corruption interpolates between correct ($\alpha{=}0$)
and model-swapped ($\alpha{=}1$) priors:
%
\begin{equation}
    \mathbf{b}_{\text{new}}^{(i)} = (1{-}\alpha)\,\mathbf{b}_{\text{orig}}^{(i)}
    + \alpha\,\mathbf{b}_{\text{orig}}^{(j)}, \quad i \neq j
\end{equation}
%
At $\alpha{=}0$, priors are correct; at $\alpha{=}0.5$, priors are
uninformative (identical for both models); at $\alpha{=}1.0$, priors
are adversarial (systematically wrong).

All experiments use $N{=}20$ independent seeds, learning rate
$\eta{=}\sqrt{\ln 2/750}{\approx}0.030$, $\gamma{=}0.05$, and
initial warmup weight $w_0{=}0.7$
(Appendix~\ref{app:initial_weight}).

% ----------------------------------------------------------------------------
% Main Result: Strategy Crossover
% ----------------------------------------------------------------------------
\paragraph{The Strategy Crossover.}
Figure~\ref{fig:prior_degradation}A presents the core finding:
\emph{strategy dominance depends on prior quality}, with a crossover
near $\alpha{\approx}$\CrossoverWC{}.

With correct priors ($\alpha{=}0$):
\begin{itemize}[nosep]
    \item \textbf{Warmup-only:} 30.0 $\pm$ 2.0 cumulative regret
    \item \textbf{Corralling:} 44.1 $\pm$ 4.7 regret (+47\% overhead)
    \item \textbf{Tabula Rasa:} 49.2 $\pm$ 3.4 regret
\end{itemize}

As priors degrade, warmup-only regret increases sharply (it has no
mechanism to detect that its priors are wrong), while Corralling's
regret increases moderately (the meta-learner shifts trust to the
tabula rasa expert).  The warmup-only and Corralling curves cross near
$\alpha{\approx}$\CrossoverWC{} (Wilcoxon signed-rank, $p{<}0.05$).
Beyond this point, warmup-only regret climbs to 98+ while Corralling
plateaus at ${\sim}67$.

Tabula rasa regret is constant at ${\sim}49$ across all $\alpha$
(by construction, it does not use priors), providing the
prior-independent baseline.

% ----------------------------------------------------------------------------
% Corralling's True Value: Variance Reduction
% ----------------------------------------------------------------------------
\paragraph{Corralling's True Value: Bounded Regret Under Uncertainty.}
Corralling's distinguishing property is
\emph{variance reduction across the unknown}.
Regret ranges across the full corruption spectrum:

\begin{itemize}[nosep]
    \item Warmup-only: 27--99 (range: 71.4)
    \item Corralling ($w_0{=}0.7$): 40--67 (range: 27.2)
    \item Tabula rasa: 49.2 (range: 0)
\end{itemize}

When a practitioner genuinely does not know their prior quality---the
most common deployment scenario---Corralling provides a bounded
worst-case guarantee.  Its maximum regret (67.3 at $\alpha{=}1$) is
32\% lower than warmup-only's worst case (98.6), while its best case
(40.5 at $\alpha{=}0.3$) is 18\% better than tabula rasa.

Crucially, this trade-off is \textbf{user-configurable} via the
\texttt{initial\_weights} parameter.
Setting $w_0{=}0.7$ (our recommended default) reduces the overhead at
$\alpha{=}0$ from 60\% (with uniform $w_0{=}0.5$) to 47\%, at the
cost of higher regret under adversarial priors (67 vs 58).
Practitioners who are more confident in their priors can increase
$w_0$ toward 0.9 to further close the gap with warmup-only, while
those seeking maximum safety can keep $w_0{=}0.5$.
Appendix~\ref{app:initial_weight} presents the full Pareto frontier
across $w_0 \in \{0.5, 0.6, 0.7, 0.8, 0.9\}$, showing that
$w_0{=}0.6$ provides a 6\% improvement at zero statistical cost
($p{=}0.49$ at $\alpha{=}1$), while $w_0 \geq 0.8$ incurs
superlinear worst-case penalties.

% ----------------------------------------------------------------------------
% Adaptive Weights
% ----------------------------------------------------------------------------
\paragraph{The Meta-Learner Correctly Adapts.}
Figure~\ref{fig:prior_degradation}B validates that Corralling's Exp4
mechanism responds correctly to prior quality across the full spectrum:

\begin{itemize}[nosep]
    \item \textbf{Strong priors ($\alpha{=}0$):} Warmup weight converges
          to ${\sim}90\%$ (correctly trusts priors)
    \item \textbf{Uninformative priors ($\alpha{=}0.5$):} Weights
          stabilise near ${\sim}75\%$ (some uncertainty)
    \item \textbf{Adversarial priors ($\alpha{=}1.0$):} Warmup weight
          drops to ${\sim}55\%$ (partially distrusts harmful priors)
\end{itemize}

Starting from $w_0{=}0.7$ rather than 0.5 means the meta-learner
reaches the correct allocation faster when priors are strong.
Under adversarial priors, the meta-learner must overcome the initial
bias, so convergence to the tabula rasa expert is slower---but it
still correctly detects and partially corrects for prior damage.

% ----------------------------------------------------------------------------
% Why the overhead exists at α=0
% ----------------------------------------------------------------------------
\paragraph{Why the Overhead Exists.}
At $\alpha{=}0$ (strong priors), the 47\% overhead arises from two
sources:
\begin{enumerate}[nosep]
    \item \textbf{Exploration cost:} Even with $w_0{=}0.7$, Corralling
          allocates ${\sim}20\%$ of queries to the tabula rasa expert
          to verify that the warmup expert is indeed superior.
    \item \textbf{Sparse signal:} 73\% of prompts produce identical
          rewards for both models, providing zero differential signal
          to the meta-learner.
\end{enumerate}

The meta-learner's Exp4 update uses importance-weighted loss
$\hat\ell_i = (1{-}r)/p_i$, which is zero whenever the selected model
succeeds.  With a theoretically optimal $\eta{\approx}0.030$,
convergence is smooth but slow---exactly as theory predicts for a
750-round horizon.

% ----------------------------------------------------------------------------
% Connection to library
% ----------------------------------------------------------------------------
\paragraph{Implications for the Library.}
This experiment validates the library's design philosophy: offering
three first-class strategies rather than a single ``best'' approach.
The crossover point provides actionable guidance:

\begin{itemize}[nosep]
    \item \textbf{Validated priors} (accuracy ${>}80\%$, $\alpha \ll$
          \CrossoverWC{}): Use warmup-only for 47\% lower regret.
    \item \textbf{Unknown prior quality} (the common case):
          Use Corralling ($w_0{=}0.7$) for bounded worst-case regret
          (max 67 vs warmup-only's 99).
    \item \textbf{Known bad or absent priors} ($\alpha \gg$
          \CrossoverWC{}): Use tabula rasa---its constant 49 regret
          is the safest option when priors are harmful.
\end{itemize}

Sections~\ref{sec:catastrophic} and~\ref{sec:zero_shot} demonstrate
additional scenarios where Corralling's adaptive capability provides
value beyond the static-benchmark analysis presented here.
